\documentclass[11pt,a4paper]{article}
\usepackage[margin=25mm]{geometry}
\usepackage{fontspec}
\setmainfont{Latin Modern Roman}
\setlength{\parindent}{0pt}
\setlength{\parskip}{0.55em}
\title{Mapa natal demonstrativo\\Síntese de abertura}
\author{Quarto lote para revisão editorial}
\date{}
\newcommand{\leitura}[3]{\subsection{#1}\textit{#2}\par #3}
\newcommand{\subtituloleitura}[1]{\textbf{#1}\par\nopagebreak[4]}
\begin{document}
\maketitle

\section*{Síntese do seu mapa}
\subtituloleitura{Você dá forma ao que percebe e ganha espaço quando põe sua própria voz na conversa}

Em um grupo, você costuma perceber cedo quem ainda não conseguiu falar e onde a conversa perdeu o fio. É natural abrir espaço com uma pergunta e, depois de ouvir, dar ao assunto uma forma que permita avançar. Seu Ascendente em Libra e seu Sol em Capricórnio na casa 3 aparecem juntos nesse gesto. Você lê a situação e devolve algo que as pessoas conseguem usar: uma explicação clara, uma decisão bem formulada ou uma maneira mais prática de começar. É assim que sua voz ganha peso com o tempo.

A Lua em Libra na casa 12 mostra o que acontece depois. Você pode conduzir bem uma conversa e só entender em silêncio o quanto ela mexeu com você. A quadratura entre Sol e Lua dá relevo a essa diferença de ritmo. Se uma decisão foi tomada antes que sua própria reação ficasse clara, você pode voltar ao assunto. Nomear o ponto que continuou incomodando costuma ser mais produtivo do que tentar acomodá-lo sozinho depois do combinado.

\subsection*{Iniciativa, compromisso e autoria}

Você enxerga possibilidades com rapidez e leva a sério o que promete. A tensão aparece quando um projeto seu também exige tempo, dinheiro ou disponibilidade de outra pessoa. Júpiter na casa 1 em oposição a Saturno na casa 7 descreve esse ponto de negociação. Lilith na casa 10, em quadratura com os dois, acrescenta a pergunta sobre autoria: quando o trabalho ganha visibilidade, fica claro o que você propôs, o que assumiu e o que foi construído em conjunto? Os três formam uma quadratura em T.

Num projeto em parceria, você pode encaminhar essa tensão logo no começo. Proponha uma primeira etapa que os dois consigam sustentar e conversem sobre como a contribuição de cada pessoa será apresentada. O combinado dá espaço para sua iniciativa continuar depois do entusiasmo inicial e torna mais fácil falar de reconhecimento quando o trabalho chegar ao público.

\subsection*{Sua direção aparece no jeito de trabalhar}

Há pistas sobre o que você quer fazer fora de casa nas tarefas que melhoram quando passam pelas suas mãos. Com o trígono entre o Meio do Céu e o Nodo Norte na casa 6, vale reparar nos ajustes que facilitam o trabalho de quem está por perto. Você pode descobrir uma direção profissional ao organizar a chegada de alguém, redesenhar uma rotina difícil ou dar forma a um cuidado que até então dependia de improviso.

O trígono entre Lilith e Quíron aproxima esse trabalho cotidiano da sua maneira de assumir autoria. Quando apresenta algo que fez, contar como começou e o que precisou ajustar dá substância à sua fala. A experiência que você ganhou testando e corrigindo a ideia se torna algo que outra pessoa consegue reconhecer, usar e talvez adaptar. É um modo de ocupar espaço público que tem a ver com o que você já faz bem no dia a dia.

\medskip
{\small\textit{Referências para continuar a leitura.} Sol em Capricórnio na casa 3, Lua em Libra na casa 12, eixo Ascendente--Descendente, casas 1, 6, 7 e 10 e os aspectos citados nesta síntese.}

\end{document}